\section{Testing and Experimentation}
\label{sec:testing}

\subsection{Methodology}
\label{subsec:methodology}

The purpose of the testing is to evaluate the performance characteristics of our custom engine in comparison with industry-standard
ECS frameworks under comparable workloads.
We assess performance through measuring system frame duration, the time required to run the user system,
as this duration is a direct indicator of runtime responsiveness and execution efficiency.

For each framework, system frame duration is measured by recording a timestamp immediately before the execution of
the user system and another immediately after.
We repeat this process for different entity counts to observe underlying performance degradation.
To evaluate the framework performance, we assume an instantaneous, highly optimized renderer.
To this end, we minimize or exclude all rendering overhead present in each framework as much as possible, where applicable.
This lets us take the system frame duration to be representative of the response time from input to output,
making it an effective metric for overall system performance.

All ECS frameworks were tested using the language’s debug build profile.
Systems for each framework were implemented according to official tutorials and documentation
to ensure that implementations were idiomatic and representative of recommended usage patterns.

All tests were executed on local hardware using the same machine and operating environment to ensure consistency across measurements.
The test system specifications are as follows:

\begin{itemize}
    \item CPU: 11th Gen Intel(R) Core(TM) i7-11800H @ 2.30 GHz
    \item GPU: Intel(R) UHD Graphics
    \item Memory: 3200 MHz 16.0 GB RAM
    \item Operating System: Windows 10 Version 22H2 (OS Build 19045.6466)
\end{itemize}

The user system used for the benchmark is a simple process that adds an entity's velocity to its position.
This system is used as it measures the framework's efficiency in iteration and data modification.

\subsubsection*{Bevy ECS}

For the Bevy ECS implementation, timing measurements are performed using the Rust standard library's \texttt{std::time::Instant}.
A benchmarking resource is defined to store frame timing information and maintain a handle to a log file for persistent output.

Two benchmark systems are inserted into the Bevy update schedule: one system records the start timestamp immediately
before execution of the user-defined systems, while another records the end timestamp immediately afterward.
Bevy’s system chaining mechanism is used to guarantee deterministic execution order between the
benchmark systems and the systems under test.
Each user system is allowed to be parallelized by Bevy's ECS scheduler.
To idiomatically measure only user systems and minimize additional overhead, we use Bevy's \texttt{MinimalPlugins} plugin group.
According to related documentation~\cite{bevy_minimalplugins}, this plugin group represents the absolute minimum bare-bones application.

The elapsed duration between the two timestamps was written both to standard output and to a log file
using \texttt{std::fs::File} guarded by a mutex for safe access.

\subsubsection*{EnTT}

For the EnTT implementation, timing measurements were performed using C++’s \texttt{<chrono>} library.
The benchmark instrumentation was placed directly within the main \texttt{update()} function responsible for executing all ECS systems.
For the user system, we use a range-for implementation as documented in accordance to the framework's given code example~\cite{ValtoLibraries_EnTT}.

A timestamp was recorded immediately before system execution using \texttt{std::chrono::steady\_clock::now()},
and a second timestamp was recorded immediately after all systems completed execution.
The elapsed duration was converted to microseconds using \texttt{std::chrono::duration\_cast}.

Measured frame times were printed to standard output and appended to a persistent log file using \texttt{std::ofstream}.

\subsubsection*{Custom Engine}

For the custom engine, frame timing was measured using the Windows high-resolution timing API \texttt{QueryPerformanceCounter}.
Timestamps were recorded before and after processing inputs and updating the scene within the main game loop.

Timing values were converted to microseconds using the system performance frequency obtained from \texttt{QueryPerformanceFrequency}.
Results were written to the engine logging system and manually extracted for later analysis and processing.